\documentclass[11pt]{article}

\usepackage[utf8]{inputenc}
\usepackage[T1]{fontenc}
\usepackage{amsmath,amssymb,amsthm,mathtools}
\usepackage{geometry}
\usepackage{graphicx}
\usepackage{booktabs}
\usepackage[round]{natbib}
\usepackage{hyperref}
\usepackage{cleveref}
\usepackage{enumitem}
\usepackage{bm}

\geometry{margin=1in}

\newtheorem{theorem}{Theorem}[section]
\newtheorem{proposition}[theorem]{Proposition}
\newtheorem{lemma}[theorem]{Lemma}
\newtheorem{corollary}[theorem]{Corollary}
\newtheorem{definition}[theorem]{Definition}
\newtheorem{assumption}[theorem]{Assumption}
\newtheorem{example}[theorem]{Example}
\theoremstyle{remark}
\newtheorem{remark}[theorem]{Remark}

\newcommand{\R}{\mathbb{R}}
\newcommand{\E}{\mathbb{E}}
\newcommand{\Prob}{\mathbb{P}}
\newcommand{\F}{\mathcal{F}}
\newcommand{\Wass}{W_2}
\DeclareMathOperator*{\argmax}{arg\,max}

\title{A Combined Particle/Finite-Difference Scheme for\\Common-Noise Mean-Field Game FBSDEs:\\Convergence Rate and CFL Condition}

\author{Nisong Monyimba\thanks{Department of Mathematics and Statistics. Correspondence: nisongmonyimba278@gmail.com. Code and supplementary material: \url{https://github.com/NisongMonyimba/AppliedMathThesis}.}}

\date{\today}

\begin{document}
\maketitle

\begin{abstract}
We develop and analyse a fully discrete numerical scheme for the conditional
McKean--Vlasov forward-backward stochastic differential equation (FBSDE) system
characterising mean-field game (MFG) equilibria with common noise. The scheme
combines an Euler--Maruyama discretisation of the forward state and backward
adjoint processes with a particle (or finite-difference) approximation of the
stochastic Fokker--Planck equation governing the population measure. We prove that
the resulting discrete equilibrium converges to the true equilibrium at rate
$O(\Delta t^{1/2} + \Delta x)$ in the Wasserstein-2 distance, and derive a
CFL-type stability condition for the coupled system that accounts separately for
parabolic diffusion, stochastic transport from the common noise, and the explicit
feedback coupling between the forward and backward components. We validate the
scheme's strong-order convergence on a tractable linear-quadratic benchmark and
discuss the practical particle-count requirements implied by the CFL condition in
higher dimensions.
\end{abstract}

\noindent \textbf{Keywords:} mean-field games, common noise, numerical methods,
particle methods, CFL condition, McKean--Vlasov FBSDEs

\noindent \textbf{MSC2020:} 65C30, 65M12, 91A16, 60H35


\section{Introduction}

The probabilistic characterisation of mean-field game (MFG) equilibria with common
noise as solutions to a conditional McKean--Vlasov forward-backward stochastic
differential equation (FBSDE) system is now well established
\citep{carmonadelarue2018,ahujarenyang2019}, and a companion paper establishes an
explicit, parameter-dependent contraction estimate for the associated mean-field
fixed-point map \citep{paper1}. Computing the equilibrium in practice, however,
requires discretising three coupled components simultaneously: the forward state
SDE, the backward adjoint BSDE, and the stochastic partial differential equation
(SPDE) governing the evolution of the population measure under the common shock.
The third component is the distinguishing numerical difficulty relative to the
classical (no-common-noise) MFG setting, where the population measure is
deterministic and solves an ordinary (not stochastic) Fokker--Planck PDE.

\paragraph{Contribution.} We propose a combined Euler--Maruyama/particle scheme for
the full coupled system (\Cref{sec:scheme}) and prove (\Cref{thm:main}) that it
converges to the true equilibrium measure flow at rate
$O(\Delta t^{1/2}+\Delta x)$ in Wasserstein-2 distance, under the same structural
assumptions used to establish well-posedness in \citet{paper1}. The proof
(\Cref{sec:proof}) combines weak-error estimates for the forward/backward
components with a strong-error estimate for the stochastic Fokker--Planck
discretisation, then closes the loop via the contraction property of the
continuous-time fixed-point map. We further derive an explicit CFL-type condition
(\Cref{sec:cfl}) for the coupled system, which -- unlike the classical parabolic CFL
condition for deterministic Fokker--Planck PDEs -- includes an additional
constraint from the stochastic transport induced by the common noise.

\paragraph{Numerical validation, and an important caveat.} We validate the scheme's
\emph{strong}-order convergence (pathwise $L^2$ error against a fine-resolution
reference path, using common random numbers) on the tractable linear-quadratic
(LQ) benchmark in \Cref{sec:experiments}, confirming the expected rate of $1.0$ for
this benchmark (additive, state-independent diffusion coefficients give the
Milstein correction term a zero coefficient, so Euler--Maruyama already attains
its maximal strong order). \textbf{We note explicitly that a full numerical
verification of the theorem's \emph{weak} (Wasserstein) rate on the LQ benchmark,
using a properly noise-controlled (common-random-number or paired) Monte Carlo
design at adequate particle counts, remains to be completed and is not reported
here; our own preliminary attempt at a naive (non-paired) weak-error sweep did not
exhibit a clean convergence trend at the particle counts we tested, indicating that
Monte Carlo sampling noise dominates the discretisation signal at those
settings.} We report only the strong-error results, which we have validated
repeatedly and which directly confirm the scheme's stability and the temporal
discretisation order; the weak-error LQ-benchmark table given in an earlier
preprint version of this material was generated using an erroneous sign in the
benchmark's closed-form Riccati coefficients and is superseded by the corrected
system given in \Cref{app:lq-corrected}.

\paragraph{Related work.} \citet{chassagneuxcrisandelarue2019} establish
convergence of a particle method for McKean--Vlasov FBSDEs without common noise.
\citet{achdouporretta2016} analyse finite-difference schemes for deterministic MFG
PDEs. \citet{lordpowellshardlow2014} provide the general SPDE numerics background
underlying our Fokker--Planck discretisation. Our contribution is the extension of
the particle/FBSDE convergence theory to the common-noise setting, where the
measure-valued component is itself stochastic, together with the explicit CFL
condition this introduces.

\paragraph{Organisation.} \Cref{sec:setup} recalls the FBSDE system.
\Cref{sec:scheme} defines the discretisation. \Cref{sec:main} states the main
convergence theorem. \Cref{sec:proof} sketches the proof. \Cref{sec:cfl} derives
the CFL condition. \Cref{sec:experiments} reports the validated strong-error
experiment. \Cref{sec:discussion} discusses scope and the open weak-error
verification.


\section{The FBSDE System}
\label{sec:setup}

We adopt the setting of \citet{paper1}: a representative agent's state
$X_t \in \R^d$ solves $dX_t = b(t,X_t,\mu_t,\alpha_t)dt + \sigma\,dW_t^i +
\sigma_0\,dW_t^0$, with $\mu_t = \mathcal{L}(X_t\mid\F_t^{W^0})$, optimal control
characterised by the stochastic maximum principle via the adjoint BSDE
$-dp_t = \partial_x H\,dt - q_t\,dW_t^i - r_t\,dW_t^0$, and the equilibrium measure
flow $\mu^*$ the unique fixed point of the mean-field map $\Phi$. Under the
standing assumptions of \citet{paper1}, $\mu^*$ satisfies the (weak form of the)
stochastic Fokker--Planck equation
\begin{equation}
    d\langle \mu_t, \phi\rangle = \langle \mu_t, \mathcal{L}_t^*\phi\rangle\, dt +
    \sigma_0\langle \mu_t, \nabla\phi\rangle\, dW_t^0, \qquad \phi \in
    C_c^\infty(\R^d),
    \label{eq:fp-spde}
\end{equation}
where $\mathcal{L}_t^*\phi = -\nabla\cdot(b\phi) + \tfrac{\sigma^2}{2}\Delta\phi$.


\section{The Discrete Scheme}
\label{sec:scheme}

Fix a time grid $t_n = n\Delta t$, $n=0,\dots,N$, $\Delta t = T/N$. The forward
state and backward adjoint are discretised by Euler--Maruyama,
\begin{align}
    X^{n+1} &= X^n + b(t_n,X^n,\mu^n,\alpha^n)\Delta t + \sigma\Delta W_n^i +
        \sigma_0\Delta W_n^0, \label{eq:disc-forward} \\
    p^n &= \E_n\big[p^{n+1} + \partial_x H(t_n,X^n,\mu^n,p^n,\alpha^n)\Delta
        t\big], \label{eq:disc-backward}
\end{align}
with the conditional expectation in \eqref{eq:disc-backward} approximated by
regression on a finite-dimensional basis \citep{gobetlemorwarin2005}. The measure
$\mu^n$ is approximated either by the empirical measure of $M$ simulated particles
or by an explicit finite-difference scheme on a spatial grid of mesh $\Delta x$ for
the strong form of \eqref{eq:fp-spde}; in the particle case there is no explicit
$\Delta x$, and the effective resolution scales as $O(M^{-1/d})$.


\section{Main Theorem}
\label{sec:main}

\begin{theorem}[Numerical convergence]
\label{thm:main}
Under the standing assumptions of \citet{paper1} and the additional smoothness
condition that $b,f,g$ are of class $C^{2,4}$ with bounded derivatives, there
exists $C>0$, independent of $\Delta t,\Delta x$, such that for all sufficiently
small $\Delta t, \Delta x$ satisfying the CFL condition \eqref{eq:cfl},
\begin{equation}
    \sup_{0\leq n\leq N} \E\big[\Wass(\mu^{\Delta,n},\mu_{t_n}^*)^2\big]^{1/2} \leq
    C(\Delta t^{1/2}+\Delta x).
    \label{eq:main-rate}
\end{equation}
\end{theorem}

\begin{proof}[Proof sketch]
See \Cref{sec:proof}.
\end{proof}


\section{Proof Sketch}
\label{sec:proof}

The proof combines three error contributions and closes the loop via the
contraction property of $\Phi$ established in \citet{paper1}.

\textbf{(i) Forward/backward weak error.} The Euler--Maruyama weak error for
Lipschitz SDEs under the stated smoothness is $O(\Delta t)$
\citep[Thm~14.5.2]{kloedenplaten1992}; an analogous $O(\Delta t)$ weak error holds
for the BSDE discretisation \citep{bouchardtouzi2004,gobetlemorwarin2005}. Both
propagate to an $O(\Delta t^{1/2})$ contribution to the Wasserstein distance between
discrete and continuous laws via the synchronous-coupling argument of
\citet{paper1}.

\textbf{(ii) Fokker--Planck discretisation error.} The explicit finite-difference
(or particle) scheme for \eqref{eq:fp-spde} satisfies, under the CFL condition
\eqref{eq:cfl} and by a standard Lax--Richtmyer consistency-plus-stability
argument adapted to the stochastic transport term \citep[Thm~10.7]{lordpowellshardlow2014},
\[
    \sup_n \E\big[\|\rho^n-\rho_{t_n}\|_{L^2}^2\big]^{1/2} \leq C(\Delta t^{1/2}+
    \Delta x).
\]

\textbf{(iii) Closing the fixed-point loop.} Let $\Phi^\Delta$ denote the discrete
fixed-point map induced by \eqref{eq:disc-forward}--\eqref{eq:disc-backward}
together with the Fokker--Planck discretisation. Combining (i)-(ii) gives
$\|\Phi^\Delta(\mu)-\Phi(\mu)\|_{\lambda,T} \leq C_{\mathrm{disc}}(\Delta
t^{1/2}+\Delta x)$ uniformly over bounded sets of measure flows. Writing
$\mu^\Delta = \Phi^\Delta(\mu^\Delta)$ for the discrete fixed point and using the
triangle inequality with the contraction factor $\rho<1$ from \citet{paper1},
\[
    \|\mu^\Delta-\mu^*\|_{\lambda,T} \leq C_{\mathrm{disc}}(\Delta t^{1/2}+\Delta x)
    + \rho\|\mu^\Delta-\mu^*\|_{\lambda,T},
\]
and rearranging gives \eqref{eq:main-rate} with
$C = C_{\mathrm{disc}}/(1-\rho)$.


\section{CFL Condition for the Coupled System}
\label{sec:cfl}

Stability of the finite-difference Fokker--Planck discretisation requires the
parabolic condition $\Delta t \leq \Delta x^2/(2\sigma^2)$. The stochastic
transport term introduced by the common noise adds the further constraint
$\Delta t \leq \Delta x/(2\|\sigma_0\|_F L_b)$ \citep[Ch.~10]{lordpowellshardlow2014}.
The explicit coupling between the discrete forward state and backward adjoint
(via $\partial_x H$ evaluated at the current discrete state) introduces a third
constraint, $\Delta t \leq 1/L_H$ with $L_H$ the Lipschitz constant of
$\partial_x H$. Collecting these,
\begin{equation}
    \Delta t \leq \min\left(\frac{\Delta x^2}{2\sigma^2}, \;
    \frac{\Delta x}{2\|\sigma_0\|_F L_b}, \; \frac{1}{L_H}\right).
    \label{eq:cfl}
\end{equation}
In the particle method, $\Delta x$ has no direct analogue; the effective
resolution scales as $O(M^{-1/d})$, so \eqref{eq:cfl} translates into a joint
requirement on the particle count $M$ and the time step that becomes increasingly
demanding as the state dimension $d$ grows -- the familiar particle-method curse of
dimensionality, here interacting with the common-noise stochastic-transport
constraint specifically.


\section{Numerical Validation: Strong-Order Convergence}
\label{sec:experiments}

We validate the temporal discretisation order on the one-dimensional
linear-quadratic benchmark with $\kappa=0.5,\lambda=1.0,\sigma=0.3,\sigma_0=0.2$,
$T=1$, comparing trajectories at four step sizes against a fine-resolution
($\Delta t = 2^{-10}$) reference path, using common random numbers (CRN) to
isolate the discretisation error from Monte Carlo sampling noise. We report the
\emph{strong} (pathwise $L^2$) error, $\E[|X^{\Delta}_T - X_T^{\mathrm{ref}}|^2]^{1/2}$,
averaged over 100 independent trials.

\begin{table}[ht]
\centering
\caption{Strong-error convergence, LQ-MFG benchmark (CRN, 100 trials)}
\label{tab:strong-error}
\begin{tabular}{lccc}
\toprule
$\Delta t$ & Mean error & Std.\ error & Estimated rate \\
\midrule
$2^{-5} = 0.03125$ & 0.01751 & 0.01001 & -- \\
$2^{-6} = 0.01563$ & 0.01030 & 0.00796 & -- \\
$2^{-7} = 0.00781$ & 0.00540 & 0.00441 & -- \\
$2^{-8} = 0.00391$ & 0.00218 & 0.00127 & \textbf{0.996} \\
\bottomrule
\end{tabular}
\end{table}

The estimated rate of $0.996$ matches the theoretical strong order for additive
(state- and measure-independent) diffusion coefficients, for which the Milstein
correction term vanishes and Euler--Maruyama already attains strong order $1$
rather than the generic order $1/2$. \textbf{This confirms the temporal
discretisation is implemented correctly and is stable, but it validates the
\emph{strong}-error behaviour of the forward SDE component in isolation; it does
not, by itself, validate the \emph{weak} (Wasserstein, measure-level) rate of
\Cref{thm:main}, which governs the full coupled scheme's convergence to the
equilibrium measure and is the theorem's actual content.} We regard a properly
noise-controlled empirical verification of the weak rate as a required next step
before this manuscript is finalised for submission (see \Cref{sec:discussion}).


\section{Discussion}
\label{sec:discussion}

\subsection{What is, and is not, empirically confirmed here}

The strong-error experiment of \Cref{sec:experiments} confirms that the forward
discretisation is implemented correctly and behaves as theory predicts for this
benchmark's specific (additive-noise) structure. It does \emph{not} confirm the
weak/Wasserstein rate $O(\Delta t^{1/2}+\Delta x)$ of \Cref{thm:main}, which is a
statement about convergence of the \emph{measure}, not a single trajectory, and
requires a separate experimental design (e.g., comparing the empirical measure of
many particles against the known closed-form Gaussian equilibrium measure via the
Wasserstein-2 formula for Gaussians, using either a much larger particle count or
a common-random-number / paired-trajectory design to suppress Monte Carlo
sampling noise enough to resolve the $O(\Delta t^{1/2})$ trend). A preliminary,
non-paired attempt at this experiment did not show a clean trend at the particle
counts tested, indicating the sampling-noise floor was not adequately controlled;
we have not yet completed a design that resolves this and do not report specific
weak-error numbers as a result.

\subsection{Practical implications of the CFL condition}

\Cref{eq:cfl}'s three terms have different practical bite depending on the
application. For low-dimensional financial models ($d \lesssim 5$), the parabolic
and coupling terms typically dominate, and standard grid resolutions suffice. For
higher-dimensional models, the particle method's $O(M^{-1/d})$ effective
resolution makes the stochastic-transport term increasingly binding, motivating
the dimension-reduction techniques (moment closures, low-rank measure
approximations) we discuss as future work.

\subsection{Limitations}

Beyond the open weak-error verification noted above: the theorem requires
$C^{2,4}$ smoothness, which excludes models with non-smooth transaction costs or
jump-type frictions common in financial applications; the particle method's
curse of dimensionality is not addressed by any variance-reduction technique in
the present scheme; and the contraction factor $\rho$ inherited from
\citet{paper1} is itself only an expectation-level (not pathwise) guarantee, which
limits the strength of any high-probability (rather than expectation) statement
one might hope to extract from \Cref{thm:main}.


\section{Conclusion}

We have proposed and analysed a combined particle/finite-difference scheme for
common-noise MFG equilibria, proving an $O(\Delta t^{1/2}+\Delta x)$ Wasserstein
convergence rate and deriving the CFL condition the coupled system requires. The
scheme's temporal discretisation is validated via a strong-error experiment on a
tractable benchmark; a corresponding empirical validation of the theorem's weak
(measure-level) rate remains open and is flagged explicitly as such, rather than
asserted without adequate experimental support.


\appendix
\section{Corrected Linear-Quadratic Benchmark Coefficients}
\label{app:lq-corrected}

For the benchmark of \Cref{sec:experiments}, the closed-loop optimal control is
$\alpha_t^* = -(A_tX_t+B_t\bar\mu_t)/\lambda$ with $(A_t,B_t)$ solving
\begin{align}
    \dot A_t &= 2\kappa A_t - A_t^2/\lambda + 2Q, & A_T&=2P, \\
    \dot B_t &= \kappa B_t + A_tB_t/\lambda - \kappa A_t, & B_T&=-2P.
\end{align}
\textbf{We flag explicitly that the sign of the cross-coupling term in the
$B_t$ equation (here, $+A_tB_t/\lambda$) was found, during preparation of this
manuscript, to have been given with the opposite sign in an earlier internal
draft; the corrected sign above was verified by an independent empirical
optimality check (perturbing the closed-loop control away from its
Riccati-implied value and confirming that realised cost strictly increases, as
required at a true optimum) before being adopted here.} For the stated parameters,
$A_t \equiv 2$ for all $t$ (a fixed point of the $A$-equation), and $B_t$ ranges
from approximately $0.20$ at $t=0$ to $-2$ at $t=T$.

\bibliographystyle{plainnat}
\bibliography{references}

\end{document}
